% Formal closure of the single post-review10 Brauer-zero main reduction.
% No separate movement event is claimed for these transfers.
\begin{corollary}[Proved: Brauer-zero counts through a detecting surface]
\label{D:Brauer-zero-detecting-surface-transfer}
Let $(S,T)$ satisfy the normal proper hypotheses of
Theorem~\ref{D:generic-surface-controlled-boundary}, in characteristic
$p>0$, and put $U=S\setminus T$ and $Q=Q(S,T)$.
Fix the surface $(X,T_X)$ and the geometric morphism
$h:X\setminus T_X\to U$ supplied there. There is a constant $C$
such that, for $q=p^f$ and $r\ge1$,
\[
 \#\left(\left\{\rho:Q\to\operatorname{PGL}_r(\mathbf F_q):
       \operatorname{br}(h^*\rho)=0\right\}
       /\operatorname{PGL}_r(\mathbf F_q)\right)
                \le q^{Cr^2}.                         \tag{D.BT.1}
\]
In particular this bounds the representations with
$\operatorname{br}(\rho)=0$ on $U$, without any bound on the
scalar centre. The same conclusion holds for every normal proper
scheme with a finite omitted subset and dimension at least two,
using Theorem~\ref{D:generic-surface-finite-punctures}; in that case
$Q=\pi_1(U)$. The chosen detecting surface is independent of $q,r$
and $\rho$.
\end{corollary}
\begin{proof}
The cited geometric theorem gives a continuous surjection
$Q(X,T_X)\twoheadrightarrow Q$. Precomposition is injective on
conjugacy classes: an equality after a surjective precomposition
is the same equality on $Q$. The representations counted on the
left of (D.BT.1) therefore inject into the Brauer-zero classes
counted in Theorem~\ref{D:Brauer-zero-projective-representation-growth}
for the one fixed normal proper surface $X$. That theorem gives
the bound with a constant depending on $(X,T_X)$, hence on the
fixed choice made from $(S,T)$.

The algebraic projective torsor and its lifting gerbe commute with
pullback, so the geometric Brauer obstruction does also. Thus
$\operatorname{br}(\rho)=0$ implies
$\operatorname{br}(h^*\rho)=0$. The converse is not used or asserted.
For a finite omitted subset, the cited finite-puncture theorem
provides precisely the same surjection of full fundamental groups.
No injectivity in degree-two cohomology is inferred from these
group surjections.
\end{proof}

\begin{corollary}[Proved: Brauer-zero counts under finite normalization]
\label{D:Brauer-zero-normalization-transfer}
Let $S$ be proper over an algebraically closed field of
characteristic $p>0$, and let $T\subset S$ have codimension at
least two in every irreducible component that it meets. Let $U$
be a connected component of $S\setminus T$. Assume either that
$T$ is finite or that the normalized components satisfy the
projective-modification hypothesis of
Corollary~\ref{D:nonnormal-controlled-boundary-quantitative-transfer}.

Let $V\to U$ be the restriction of the finite normalization of
$S_{\mathrm{red}}$, with connected components $V_i$. For each
component of dimension at least two, fix one detecting surface
map $h_i:W_i\to V_i$ as in the preceding corollary; here $W_i$
is a normal proper surface with a finite set removed.
There is a constant $C$ such that the number of conjugacy classes
of representations
\[
                \rho:\pi_1(U)\to
                          \operatorname{PGL}_r(\mathbf F_q)
\]
whose geometric Brauer class vanishes after pullback to every
such $W_i$ is at most $q^{Cr^2}$, for all $q=p^f$ and $r\ge1$.
In particular this holds for all representations with zero
geometric Brauer class on $U$. No condition on the scalar-centre
order is imposed.
\end{corollary}
\begin{proof}
The reduction and normalization constructions, the codimension
claim, and the fixed detecting surfaces are exactly those in
Corollaries~\ref{D:nonnormal-finite-puncture-quantitative-transfer}
and~\ref{D:nonnormal-controlled-boundary-quantitative-transfer}.
The preceding corollary bounds the allowed restriction classes
on every $V_i$ of dimension at least two by $q^{C_i r^2}$.
Components of dimension at most one are normal and proper,
since the componentwise codimension hypothesis leaves no
omitted point on them. Their fundamental groups are finitely
generated by the already established projective cases.
If $d_i$ generators suffice, the number of homomorphisms into
$P=\operatorname{PGL}_r(\mathbf F_q)$ is at most
$|P|^{d_i}\le q^{d_i r^2}$. No Brauer restriction is needed on
these components.

Set $a=\#\pi_0(V\times_U V)$, which is finite. Fix one tuple of
allowed $P$-torsor classes on the $V_i$ and choose representative
torsors. The Isom scheme between their two pullbacks to
$V\times_U V$ is finite etale and has $|P|$ elements in each
geometric fibre. A section on a connected component is determined
by its value in one geometric stalk. Hence at most $|P|^a$
descent isomorphisms are possible. The cocycle condition and
passage to isomorphism classes only decrease that number.
Finite proper descent of finite etale schemes, including the
action morphisms, is effective and fully faithful
\citep[Expose IX, Theorem~4.7]{C_SGA1}; this is the same explicit
descent argument used in the two cited normalization corollaries.

Multiplying the component bounds and the descent bound gives
$q^{Cr^2}$ with
$C=a+\sum_{\dim V_i\ge2}C_i+\sum_{\dim V_i\le1}d_i$.
We may count every descent datum on a fixed allowed tuple, even
if its descended Brauer class on $U$ is nonzero. The resulting
overcount is valid for the stated condition, which was imposed
only after the fixed surface pullbacks. Finally zero Brauer class
on $U$ pulls back to zero on each $W_i$. Again no degree-two
cohomology injection or descent theorem is asserted.
\end{proof}

\begin{remark}[Scope and movement accounting]
These are the functorial and finite-normalization closure of the
single Brauer-zero projective-growth theorem. They do not supply
an additional independent main-gate reduction. They also do not
bound projective representations with nonzero Brauer class on
the fixed detecting surfaces, nor arbitrary nonabelian crown
multiplicities, and they do not settle general D.1 or D.2(ii).
\end{remark}
